% Copyright (c) 2026 Eugene Joseph M. Ragasa.
% SPDX-License-Identifier: Apache-2.0
% Author: Eugene Joseph M. Ragasa
% Formal proof completed: 2026-10-20
\documentclass[11pt]{article}
\usepackage[T1]{fontenc}\usepackage{amsmath,amssymb}\usepackage[margin=1in]{geometry}\usepackage[hidelinks]{hyperref}
\title{Proof Dossier: Positive-Outcome State Normalization}\author{Eugene Joseph M. Ragasa}\date{Formal proof completed: 2026-10-20}
\begin{document}\maketitle
\section{Identity and status}
ID \texttt{QM-MEASUREMENT-THM-003}; status \texttt{proved}. Lean declaration:
\par\noindent{\scriptsize\path{ProjectKoios.QuantumMechanics.Measurement.norm_postMeasurementState_vector_eq_one}}.
\section{Exact statement}
If the defined probability $\lVert M_i\psi\rVert^2$ is strictly positive, the selected updated
representative
\[\psi_i=\lVert M_i\psi\rVert^{-1}M_i\psi\]
has norm one.
\section{Source and assumptions}
This is source-neutral normalization mathematics. The strict positivity premise excludes a zero
operation output. No assertion is made about collapse as a universal physical postulate or about
measurements outside the declared finite bounded-operation structure.
\section{Proof schematic and proof}
Strict positivity implies $M_i\psi\ne0$. A nonzero vector divided by its norm has norm one. The
constructed value is stored as a \texttt{NormalizedStateVector}, whose normalization field yields
the exact conclusion.
\section{Lean realization and trusted base}
The construction uses Mathlib's \texttt{NormedSpace.normalize}; its proof rejects the zero case by
contradicting the positive Born weight. The theorem projects the resulting normalization field.
Kernel-reported axioms are \texttt{propext}, \texttt{Classical.choice}, and \texttt{Quot.sound};
there are no admissions.
\section{Boundary and review}
No outcome-selection dynamics, probability interpretation beyond the definition, experiment, or
continuous-spectrum update is established. AI-assisted and kernel-checked; no independent expert
review is recorded.
\end{document}
